% !PRINTED 41
\weberpage{41}

Wir unterscheiden vier F\"alle:

1. $f''''(\alpha)$ und $f(\alpha)$ nicht $=0$ [d.\,h. $f(x)$ vom vierten Grade und $\alpha$ keine Wurzel von $f(x)$].

In diesem Falle l\"a\ss t sich durch die Formel (2) f\"ur $n=0$ $S_3$ durch $S_0, S_1, S_2$ ausdr\"ucken, und daher $S_n$ f\"ur jedes positive $n$ durch $S_0, S_1, S_2$. Setzt man aber $n=-1,-2,\ldots$, so erh\"alt man $S_{-2}, S_{-3}, \ldots$ ausgedr\"uckt durch $S_{-1}, S_0, S_1, S_2$. Also bleiben in diesem Falle die vier Grundintegrale
\[
S_{-1}, \quad S_0, \quad S_1, \quad S_2.
\]

2. $f''''(\alpha)$ nicht $=0$, $f(\alpha)=0$. In diesem Falle kann man noch $S_{-1}$ durch $S_0, S_1, S_2$ ausdr\"ucken und erh\"alt als Grundintegrale
\[
S_0, \quad S_1, \quad S_2.
\]

3. $f''''(\alpha)=0$, $f(\alpha)$ nicht $=0$. Hier sind die Grundintegrale
\[
S_{-1}, \quad S_0, \quad S_1.
\]

4. $f''''(\alpha)=0$, $f(\alpha)=0$. Hier sind die Grundintegrale
\[
S_0, \quad S_1.
\]

Das Resultat hiervon ist also:

I. Ist $f(x)$ vom vierten Grade ($a_0=1$), so lassen sich alle Integrale $\Omega(x)$ ausdr\"ucken durch Integrale der Form
\[
S_0 = \int \frac{dx}{\sqrt{f(x)}}, \quad S_1 = \int \frac{x\,dx}{\sqrt{f(x)}}, \quad S_2 = \int \frac{x^2\,dx}{\sqrt{f(x)}},
\]
\[
S_{-1} = \int \frac{dx}{(x-\alpha)\sqrt{f(x)}},
\]
worin $S_{-1}$ noch von dem Parameter $\alpha$ abh\"angig ist.

II. Ist aber $f(x)$ vom dritten Grade, so gen\"ugen die drei Integrale
\[
S_0 = \int \frac{dx}{\sqrt{f(x)}}, \quad S_1 = \int \frac{x\,dx}{\sqrt{f(x)}}, \quad S_{-1} = \int \frac{dx}{(x-\alpha)\sqrt{f(x)}}.
\]

In beiden F\"allen ist $S_0$ \gesperrt{von der ersten Gattung}. Ist $f(x)$ vom dritten Grade, so ist $S_1$ von der zweiten und $S_{-1}$ von der dritten Gattung. Ist dagegen $f(x)$ vom vierten Grade, so sind $S_1$, $S_2$ und $S_{-1}$ von der dritten Gattung. Man kann aber aus $S_1$ und $S_2$ ein Integral zweiter Gattung zusammensetzen. Denn es ist
%CONT>>
